% Compile with pdfLaTeX on Overleaf
% !TeX program = pdflatex
\documentclass[conference]{IEEEtran}

\usepackage[utf8]{inputenc}
\usepackage[T5]{fontenc}
\usepackage[vietnamese,english]{babel}
\usepackage{lmodern}
\usepackage{amsmath,amssymb}
\usepackage{booktabs,array}
\usepackage{tikz}
\usetikzlibrary{arrows.meta,positioning}
\usepackage{url}
\usepackage[hidelinks]{hyperref}

\newcommand{\listen}{\ifmmode\mathrm{LISTEN}\else\textsc{Listen}\fi}
\newcommand{\commit}{\ifmmode\mathrm{COMMIT}\else\textsc{Commit}\fi}
\newcommand{\lcp}{\operatorname{LCP}}

\title{Nghiên cứu mô hình quyết định thời điểm dịch\\
cho hệ thống dịch Anh--Việt gia tăng độ trễ thấp}

\author{\IEEEauthorblockN{Đặng Quang Minh -- MSV: 24022397}
\IEEEauthorblockA{Trường Đại học Công nghệ -- Đại học Quốc gia Hà Nội\\
Mã nguồn: \href{https://github.com/MinhCYB/TimelyMT}{github.com/MinhCYB/TimelyMT}}}

\begin{document}
\maketitle

\begin{abstract}
Dịch gia tăng phải công bố đầu ra trước khi toàn bộ phát ngôn nguồn kết thúc, do đó chất lượng không chỉ phụ thuộc bộ dịch mà còn phụ thuộc thời điểm phân đoạn và chốt bản dịch. Bài báo nghiên cứu câu hỏi: khi bộ dịch Anh--Việt được giữ cố định, việc đưa lịch sử đã chốt vào bộ phân loại biên \listen/\commit{} có cải thiện đánh đổi chất lượng--độ trễ hay không? Hệ thống dùng EnViT5 cố định để dịch các prefix nguồn nhân quả và logistic regression để dự đoán biên. Ba biến thể tạo thành một history ablation có kiểm soát: P0 không dùng văn bản lịch sử, P1 bổ sung đơn vị nguồn vừa chốt, và P2 bổ sung thêm đầu ra đích do chính hệ thống sinh. Trên frozen DEV, quy tắc đăng ký trước chọn P1 tại threshold 0,60, đạt 27,09 BLEU và 61,40 chrF2. Tuy nhiên LA-2 đạt chất lượng cao hơn với độ trễ thấp hơn theo cả AL và LAAL; lịch sử nguồn hoặc đích cũng không tạo cải thiện nhất quán qua các threshold. Vì TEST held-out chưa được chạy, kết quả chỉ mô tả development evaluation và không chứng minh khả năng khái quát hóa held-out.
\end{abstract}

\begin{IEEEkeywords}
dịch gia tăng, dịch Anh--Việt, độ trễ thấp, segmentation policy, Local Agreement, history ablation
\end{IEEEkeywords}

\section{Giới thiệu}
Dịch gia tăng (incremental translation) quan sát phát ngôn nguồn theo thời gian và phải phát đầu ra trước khi toàn bộ nguồn kết thúc. Khác dịch câu ngoại tuyến, hệ thống phải quyết định đồng thời \emph{dịch gì} và \emph{khi nào công bố}. Chốt quá sớm làm giả thuyết thiếu ngữ cảnh; chờ quá lâu làm mất lợi ích tương tác thời gian thực. Vì đầu ra đã công bố được xem là bất biến, một quyết định biên sai không thể được sửa như suffix chưa hiển thị.

TimelyMT tách hai vai trò. \emph{Translator} ánh xạ prefix tiếng Anh đã quan sát thành giả thuyết tiếng Việt; \emph{policy} chỉ quyết định \listen{} hay \commit. EnViT5 được giữ cố định, không được fine-tune. Nhờ đó, thay đổi chất lượng--độ trễ có thể được quy cho policy thời điểm thay vì thay đổi năng lực mô hình dịch.

Policy độ dài/thời gian cố định dễ tái lập nhưng không phản ứng với độ ổn định nội dung. Phương pháp dựa trên agreement phản ứng với biến thiên giả thuyết nhưng dùng quy tắc cứng. Policy học có thể kết hợp tín hiệu prefix và trạng thái hiện tại. Câu hỏi nghiên cứu ở đây là: \emph{lịch sử văn bản đã commit có cung cấp tín hiệu bổ sung hữu ích cho quyết định biên nhân quả hay không?}

Các đóng góp được repository hỗ trợ gồm:
\begin{itemize}
  \item benchmark streaming Anh--Việt từ 17 bài TED kỹ thuật, với luồng token nguồn và source-clock mô phỏng được đóng băng;
  \item so sánh fixed, heuristic, literature adaptation và learned policy dưới cùng EnViT5 cố định;
  \item history ablation P0/P1/P2 để cô lập tác dụng của lịch sử nguồn đã chốt và lịch sử đích do hệ thống sinh;
  \item protocol tách oracle tạo nhãn khỏi suy luận nhân quả, cùng metric và quy tắc chọn cấu hình DEV xác định trước.
\end{itemize}
Đây là thiết kế thí nghiệm và implementation research có khả năng tái lập, không phải tuyên bố về một mô hình dịch mới. Báo cáo định lượng hiện chỉ dựa trên frozen DEV; TEST held-out chưa được chạy.

\section{Nghiên cứu liên quan}
\subsection{Phân đoạn cố định và độ trễ}
Chiến lược cố định theo số token hoặc thời gian xác định trực tiếp tần suất dịch, không cần huấn luyện, nhưng không biết prefix đã đủ ổn định hay chưa. Average Lagging (AL) đo mức chậm so với tác tử đọc--ghi lý tưởng~\cite{ma2019stacl}; LAAL điều chỉnh AL để việc sinh thừa không được tưởng thưởng~\cite{papi2022laal}. TimelyMT báo cáo cả hai trên vị trí lexical source token.

\subsection{Agreement và ổn định giả thuyết}
Local Agreement công bố prefix chung của các giả thuyết gần nhau và giữ đầu ra đã phát là bất biến. TimelyMT cài đặt adaptation LA-2 dựa trên hai giả thuyết liên tiếp, theo ý tưởng được Liu et al. dùng cho nhận dạng và dịch nói độ trễ thấp~\cite{liu2020partial}. Hai baseline \texttt{local\_agreement\_style\_k2/k3} khác LA-2: chúng dùng tỷ lệ longest common prefix (LCP) của một cửa sổ để chốt \emph{toàn bộ} unit. Vì vậy chúng là heuristic nội bộ, không phải reproduction của Local Agreement nguyên bản.

\subsection{Phân đoạn thích nghi học được}
Zhang et al. học phân đoạn thành Meaningful Units (MU), có sử dụng tín hiệu tương thích với bản dịch~\cite{zhang2020adaptive}. TimelyMT triển khai một \emph{adaptation}: nhãn được suy ra từ mức bảo toàn prefix giữa bản dịch hiện tại và bản dịch phần nguồn còn lại hợp lệ, rồi huấn luyện logistic regression. Cài đặt này giữ nguyên nguyên lý phân đoạn có nhận biết bản dịch nhưng không tái tạo kiến trúc neural, mục tiêu reinforcement, ASR, cặp ngôn ngữ hay chế độ huấn luyện NMT của công trình gốc.

\section{Phương pháp}
\subsection{Pipeline nhân quả}
Hình~\ref{fig:pipeline} trình bày vòng lặp. Tại bước $t$, runtime chỉ mở prefix đã phát. Bộ dịch tạo giả thuyết $H_t$ từ buffer chưa chốt. Policy dùng trạng thái nhân quả $z_t$ để quyết định
\begin{equation}
a_t=\begin{cases}
\commit,&p_\theta(\commit\mid z_t)\geq\eta,\\
\listen,&\text{ngược lại},
\end{cases}
\label{eq:decision}
\end{equation}
với threshold $\eta$. Khi \commit, giả thuyết được nối vào đầu ra và buffer mới bắt đầu sau biên; khi \listen, hệ thống chờ token kế tiếp. Độ dài tối thiểu là 4 token và 48 token là biên cưỡng bức. Translator EOS chỉ kết thúc một lần sinh giả thuyết, không đồng nghĩa policy \commit.

\begin{figure}[t]
\centering
\begin{tikzpicture}[
 node distance=2.5mm and 5mm,
 box/.style={draw,rounded corners,align=center,font=\scriptsize,
   text width=2.7cm,minimum height=6mm,inner sep=2mm},
 branch/.style={draw,rounded corners,align=center,font=\scriptsize,
   text width=2.45cm,minimum height=7mm,inner sep=1.5mm},
 flow/.style={-{Latex[length=1.5mm]},thick}]
\node[box] (stream) {Token nguồn EN đã phát};
\node[box,below=of stream] (buffer) {Buffer nguồn hiện tại};
\node[box,below=of buffer] (translator) {EnViT5 cố định\\$H_t=T(X_{\leq t})$};
\node[box,below=of translator] (policy) {Policy nhân quả\\$p_\theta(\commit\mid z_t)$};
\node[box,below=of policy] (action) {Quyết định\\\listen{} / \commit};
\node[branch,below=of action,xshift=-1.55cm] (listen-node) {\listen{}\\quay về buffer};
\node[branch,below=of action,xshift=1.55cm] (history) {\commit{}\\cập nhật committed history};
\draw[flow] (stream)--(buffer);
\draw[flow] (buffer)--(translator);
\draw[flow] (translator)--(policy);
\draw[flow] (policy)--(action);
\draw[flow] (action.south west)--(listen-node.north);
\draw[flow] (action.south east)--(history.north);
\draw[flow] (listen-node.west)--++(-3mm,0)|-(buffer.west);
\draw[flow] (history.east)--++(3mm,0)|-(policy.east);
\end{tikzpicture}
\caption{Translator cố định sinh giả thuyết hiện tại; policy nhân quả chọn \listen{} hoặc \commit. Nhánh \commit{} cập nhật committed history cho bước tiếp theo.}
\label{fig:pipeline}
\end{figure}

\subsection{Biểu diễn nguồn streaming}
Mỗi talk là dãy lexical token tiếng Anh $(x_1,\ldots,x_n)$ với thời điểm hoàn tất \texttt{emit\_ms}. Tokenization độc lập với tokenizer neural; whitespace và dấu câu đứng riêng bị loại, còn dấu nháy/gạch nối trong từ được giữ. Source-clock TED được mô phỏng ở 2,5 token/giây theo segment; thời lượng mặc định phân bổ theo số ký tự chữ/số. Quan sát runtime không có segment tương lai, reference punctuation tương lai hay alignment oracle. Timing này phục vụ comparison có kiểm soát, không đại diện latency acoustic/ASR thực.

\subsection{Bộ dịch cố định}
Mọi strategy gọi \texttt{VietAI/envit5-translation}, model được công bố cùng MTet~\cite{ngo2022mtet}, ở revision \texttt{840bc88104d5a4277af740eaedb024df8c3093e7}. Giải mã greedy dùng \texttt{do\_sample=false}, \texttt{num\_beams=1}, \texttt{max\_new\_tokens=256}. Thiết bị tự chọn CUDA/FP16 nếu có, nếu không dùng CPU/FP32; cache được phân tách theo device/dtype. Input tối đa 512 model token gồm prefix \texttt{en:}; quá giới hạn gây lỗi thay vì cắt ngầm.

Wrapper không hoàn thiện câu, thêm dấu câu hay trim nguồn. Sau decode, nó chỉ bỏ đúng tag \texttt{vi:} ở đầu cùng tối đa một dấu cách ASCII; không chuẩn hóa Unicode, case hoặc dấu câu. Policy và evaluator dùng giả thuyết đã chuẩn hóa này. Thiết kế cố định model/revision/decoding cô lập policy, nhưng vẫn cần ghi device/dtype trong artifact vì số học có thể khác nhau.

\subsection{Nhãn giả từ độ ổn định tương lai}
Trên TRAIN/DEV, tại prefix đủ 4 token, hệ thống dịch prefix hiện tại và hai cập nhật kế tiếp nếu cả hai nằm trong giới hạn 48 token. Với $H_t,H_{t+1},H_{t+2}$ và whitespace tokenization,
\begin{equation}
r_t=\frac{|\lcp(H_t,H_{t+1},H_{t+2})|}{\max(1,|H_t|)}.
\label{eq:stability}
\end{equation}
Nhãn là \commit{} khi $r_t\geq0{,}90$; talk end và độ dài 48 là biên cưỡng bức. Nếu không có đủ đúng hai cập nhật tương lai thì không áp dụng tiêu chí ổn định. Sau \commit, episode mới bắt đầu ở token kế tiếp.

$H_{t+1}$ và $H_{t+2}$ là \textbf{oracle supervision chỉ dùng khi tạo nhãn}. Chúng không nằm trong feature, checkpoint input hay rollout. Code từ chối pseudo-label TEST và validator kiểm tra leakage theo tên feature. Policy cuối chỉ fit từ manifest TRAIN checksum hợp lệ.

\subsection{Policy và history ablation}
Policy là logistic regression với \texttt{class\_weight=balanced}, solver \texttt{liblinear}, tối đa 1000 vòng lặp và seed 20260809. Mỗi trường text dùng TF--IDF word unigram/bigram, \texttt{min\_df=2}, tối đa 10.000 feature. Mười một feature số được chuẩn hóa: kích thước buffer theo token/ký tự và source-clock; độ dài giả thuyết hiện tại/trước và delta; LCP/change ratio giữa hai giả thuyết trong episode; số unit đã commit; kích thước unit nguồn/đích vừa commit.

\begin{table}[t]
\caption{History ablation của policy TimelyMT}
\label{tab:ablation}
\centering
\scriptsize
\begin{tabular}{p{0.10\columnwidth}p{0.20\columnwidth}p{0.25\columnwidth}p{0.30\columnwidth}}
\toprule
Biến thể & Local/causal & Nguồn đã chốt & Đích hệ thống \\
\midrule
P0 & Có & Không & Không \\
P1 & Có & Có & Không \\
P2 & Có & Có & Có \\
\bottomrule
\end{tabular}
\end{table}

P0, P1 và P2 dùng cùng feature số; ablation chỉ đổi text field (Bảng~\ref{tab:ablation}). P0 không có \emph{committed text history}, nhưng feature số vẫn chứa trạng thái cục bộ và kích thước unit trước. P1 thêm văn bản nguồn của unit vừa commit. P2 thêm văn bản đích do hệ thống tự sinh, không dùng gold target. Rollout là tuần tự: quyết định trước xác định buffer và history ở bước sau, không dùng teacher-forced history.

\subsection{Meaningful Unit adaptation}
MU có supervision độc lập với (\ref{eq:stability}). Với candidate $S_{s:t}$, oracle TRAIN/DEV là $H_{\mathrm{full}}$, bản dịch đoạn bắt đầu tại $s$ và kết thúc ở talk end hoặc token thứ 48. Biên MU được gán khi
\begin{equation}
\frac{|\lcp(H_t,H_{\mathrm{full}})|}{\max(1,|H_t|)}\geq0{,}90.
\end{equation}
Runtime không nhận $H_{\mathrm{full}}$: logistic regression threshold 0,50 chỉ dùng source buffer, giả thuyết hiện tại và năm feature số độ dài/thời gian. TF--IDF dùng unigram/bigram, tối đa 10.000 feature/trường và \texttt{min\_df=1}. TEST supervision bị cấm. Đây là literature adaptation, không phải reproduction hoàn hảo của Zhang et al.

\section{Thiết lập thực nghiệm}
\subsection{Dataset và split}
TimelyMT Streaming Dataset v1 gồm 17 bài TED Anh--Việt về AI, ML, NLP, robotics, computer science và chủ đề kỹ thuật liên quan. Pipeline gồm acquisition, parsing, căn chỉnh segment EN--VI đơn điệu xác định, timing nguồn, canonicalization và validation. Aligner cho phép nhóm N:M tối đa 4 với \texttt{skip\_penalty=1.6}, \texttt{group\_penalty=0.65}; nó không có mô hình lexical/semantic song ngữ.

\begin{table}[t]
\caption{Thống kê TimelyMT Streaming Dataset v1}
\label{tab:dataset}
\centering
\footnotesize
\begin{tabular}{lr}
\toprule
Thuộc tính & Giá trị \\
\midrule
Nhà cung cấp / hướng dịch & TED / EN$\rightarrow$VI \\
Số talk & 17 \\
Source / target segments & 2.471 / 2.372 \\
Alignment units & 2.323 \\
Runtime lexical tokens & 41.739 \\
Tổng source-clock mô phỏng & 16.695.600 ms \\
TRAIN / DEV / TEST (talk) & 12 / 3 / 2 \\
\bottomrule
\end{tabular}
\end{table}

Split ở cấp talk, nhóm theo speaker bằng \texttt{seeded\_sorted\_groups}, seed 20260809; không random sentence/prefix. Kiểm tra speaker leakage đã đạt. Ba talk calibration bị loại khỏi TEST. TEST không được dùng để tạo pseudo-label/MU supervision, thiết kế feature, tuning threshold hay chọn variant. CLI chỉ mở TEST sau khi có \texttt{frozen-eval-config.json}.

Toàn bộ 17 talk dùng timing mô phỏng do transcript TED không cung cấp subtitle timing dùng được. \texttt{emit\_ms} là thời điểm token hoàn tất. Tổng thời lượng trong Bảng~\ref{tab:dataset} là source-clock tổng hợp, không phải audio đã đo; mọi thống kê thời gian vì vậy là \emph{simulated source-clock latency}.

\subsection{Strategies và cấu hình}
Fixed-N phân đoạn mỗi 4/8/12 token. Fixed-Time chốt sau tối thiểu 4 token khi elapsed time đạt 1.600/3.200/4.800 ms, với trần 48. Local-Agreement-style dùng cửa sổ $k=2,3$ và LCP ratio 0,90 để commit toàn unit. LA-2 bắt đầu dịch từ token thứ tư, so sánh hai giả thuyết liên tiếp và chỉ phát prefix đích mới ổn định; suffix được flush ở talk end hoặc trần 48. Sự kiện LA-2 là target-prefix emission, nên thống kê unit của nó diễn giải khác whole-unit commit.

\begin{table}[t]
\caption{Policy và baseline đã cài đặt}
\label{tab:strategies}
\centering
\scriptsize
\begin{tabular}{p{0.27\columnwidth}p{0.20\columnwidth}p{0.44\columnwidth}}
\toprule
Policy & Loại & Tham số chính \\
\midrule
Fixed-N & cố định & $N\in\{4,8,12\}$ token \\
Fixed-Time & cố định & $\Delta\in\{1600,3200,4800\}$ ms \\
LA-style $k2/k3$ & heuristic & $k\in\{2,3\}$; ratio 0,90 \\
Local Agreement LA-2 & adaptation & 2 giả thuyết; phát LCP mới \\
MU Zhang-2020 & adaptation học & threshold 0,50 \\
P0/P1/P2 & history ablation & $\eta\in\{.30,.40,.50,.60,.70\}$ \\
\bottomrule
\end{tabular}
\end{table}

Mọi strategy dùng cùng frozen translator và cùng source stream. Điều này tránh confound do đổi translator, dù semantics event của LA-2 khác whole-unit policy.

\subsection{Metrics}
Prediction talk là phép nối bằng một dấu cách giữa các phần đích đã commit. Chỉ evaluator nạp reference tiếng Việt và tính corpus BLEU cùng chrF, $\beta=2$, bằng SacreBLEU~\cite{post2018sacrebleu}; version và signatures được lưu trong artifact.

Với $|X|$ token nguồn, $|Y^*|$ token hypothesis và $g(t)$ là số token nguồn đã quan sát khi token đích $t$ được phát, implementation tính
\begin{equation}
\operatorname{AL}=\frac{1}{\tau}\sum_{t=1}^{\tau}
\left[g(t)-\frac{t-1}{|Y^*|/|X|}\right].
\label{eq:al}
\end{equation}
$\tau$ là token đích đầu tiên có $g(t)\geq|X|$, hoặc $|Y^*|$ nếu không có. LAAL thay $|Y^*|$ trong tỷ lệ bằng $\max(|Y^*|,|Y|)$ với $|Y|$ là số whitespace token reference. AL/LAAL được tính theo talk rồi trung bình có trọng số theo số token đích. Đơn vị là lexical source token, không phải ms. Chỉ số hỗ trợ gồm commit/100 source token, mean/median token/unit, source-clock/unit, first-commit latency và forced commit rate; forced commit là tỷ lệ event do \texttt{max\_length} hoặc \texttt{talk\_end}.

\subsection{Chọn cấu hình DEV}
Chỉ 15 cấu hình learned P0/P1/P2 tham gia selection; MU và LA-2 chỉ để so sánh. AL của \texttt{fixed\_n\_8} là reference. Trong các candidate có AL không lớn hơn reference, chọn chrF2 cao nhất; tie-break theo BLEU cao hơn, AL thấp hơn, rồi threshold lớn hơn. Nếu không candidate nào thỏa, chọn AL nhỏ nhất rồi dùng cùng quality key. Selection từ chối smoke/partial metrics. Frozen config sau selection phải chứa checksum dữ liệu/split, translator identity, feature và pseudo-label config, checkpoint hashes, variant và threshold trước TEST.

\section{Kết quả và thảo luận}
\subsection{Frozen DEV và lựa chọn cấu hình}
Checkpoint được audit có stage \texttt{dev-frozen-complete}; toàn bộ metric trong phần này có \texttt{artifact\_status=full} và \texttt{publishable=true}. Quy tắc chọn dùng AL của Fixed-N/8 làm mốc ($-36{,}06$). Không learned candidate nào có AL không lớn hơn mốc này, nên nhánh dự phòng chọn AL nhỏ nhất rồi mới xét chất lượng. Kết quả là P1 tại $\eta=0{,}60$, đúng với \texttt{dev-selection.json} và \texttt{frozen-eval-config.json}. Các checksum dataset/split, revision translator và hash checkpoint đã được đóng băng trước khi mở TEST. TEST held-out chưa chạy và không có số TEST nào được trình bày dưới đây.

\begin{table}[t]
\caption{Kết quả frozen DEV. P0/P1/P2 dùng cùng $\eta=0{,}60$; in đậm là cấu hình được chọn. AL/LAAL tính theo token nguồn.}
\label{tab:results}
\centering
\scriptsize
\setlength{\tabcolsep}{3.2pt}
\begin{tabular}{lrrrr}
\toprule
Policy & BLEU & chrF2 & AL & LAAL \\
\midrule
Fixed-N/4 & 25,17 & 61,69 & $-8,24$ & $-8,24$ \\
Fixed-N/8 & 27,23 & 61,40 & $-36,06$ & $-9,81$ \\
Fixed-N/12 & 29,09 & 61,45 & $-28,54$ & 35,93 \\
Fixed-Time/1600 & 27,31 & 62,25 & $-1,82$ & 19,35 \\
Fixed-Time/3200 & 29,62 & 62,60 & 10,15 & 64,17 \\
Fixed-Time/4800 & 29,02 & 62,01 & 6,66 & 101,48 \\
LA-style $k=2$ & 26,78 & 61,29 & 0,64 & 61,43 \\
LA-style $k=3$ & 28,65 & 61,07 & $-2,23$ & 137,52 \\
Local Agreement LA-2 & 28,30 & 62,00 & $-9,35$ & 51,16 \\
MU Zhang-2020 & 28,87 & 61,56 & $-7,35$ & 60,39 \\
P0/0,60 & 27,47 & 61,69 & 1,71 & 60,28 \\
\textbf{P1/0,60} & \textbf{27,09} & \textbf{61,40} & \textbf{$-8,92$} & \textbf{61,25} \\
P2/0,60 & 27,88 & 62,01 & 1,10 & 56,53 \\
\bottomrule
\end{tabular}
\end{table}

\subsection{So sánh quality--latency}
Learned policy có thể cạnh tranh ở một số điểm nhưng không vượt các fixed policy nói chung. P1/0,60 có BLEU gần Fixed-N/8 (27,09 so với 27,23) nhưng chrF2 tương đương và LAAL cao hơn 71,06 token; Fixed-N/8 vì vậy là baseline đơn giản đặc biệt đáng chú ý. Fixed-Time/3200 cho chất lượng cao nhất bảng, hơn P1/0,60 2,53 BLEU và 1,20 chrF2, đổi lại AL tăng 19,07 token và LAAL tăng 2,92 token. Fixed-Time/1600 cũng hơn P1 về cả BLEU và chrF2, với LAAL thấp hơn đáng kể dù AL cao hơn 7,10 token. Các khác biệt giữa AL và LAAL cho thấy over-generation ảnh hưởng mạnh đến diễn giải latency; không nên kết luận từ riêng AL âm.

So với agreement, P1/0,60 bị LA-2 trội hơn: LA-2 tăng 1,21 BLEU và 0,60 chrF2, đồng thời giảm 0,43 AL và 10,09 LAAL. LA-style $k=3$ đạt BLEU 28,65 nhưng chrF2 thấp hơn P1 và LAAL rất cao; $k=2$ không tạo ưu thế rõ. MU Zhang-2020 tăng 1,78 BLEU và 0,16 chrF2 so với P1; AL của MU cao hơn 1,57 token nhưng LAAL thấp hơn 0,86 token. Do LA-2 phát target prefix còn các policy khác commit whole unit, đây là so sánh trade-off đầu ra, không phải bằng chứng ưu thế kiến trúc tổng quát.

\subsection{History ablation}
Ở cùng threshold 0,60, thêm nguồn đã chốt (P1 so với P0) giảm AL 10,63 token nhưng giảm 0,38 BLEU và 0,29 chrF2, trong khi LAAL tăng 0,97 token. Ở threshold 0,40 và 0,50, P1 lại cải thiện chất lượng và giảm cả AL lẫn LAAL so với P0; tại 0,30 và 0,70 kết quả tiếp tục pha trộn. Vì vậy DEV cho thấy source history có thể thay đổi điểm vận hành, nhưng không chứng minh một cải thiện bền vững qua threshold.

Thêm đích do hệ thống sinh cũng không cho lợi ích nhất quán. Tại 0,60, P2 hơn P1 0,79 BLEU và 0,61 chrF2; AL tăng 10,02 token nhưng LAAL giảm 4,72 token. So với P0 tại cùng threshold, P2 tăng 0,41 BLEU và 0,32 chrF2, đồng thời giảm AL 0,61 token và LAAL 3,75 token. Tuy nhiên thứ hạng này không ổn định trên toàn bộ năm threshold đăng ký trước, và quy tắc selection cuối vẫn chọn P1 vì AL nhỏ nhất trong learned candidates. Evidence frozen DEV do đó không đủ để khẳng định committed source history hay system-generated target history cải thiện trade-off tổng quát.

Về hành vi commit, P1/0,60 dùng trung bình 7,66 token nguồn/unit, 13,05 commit trên 100 token nguồn, first commit sau 20,33 token (8.009 ms source-clock mô phỏng), và forced-commit rate chỉ 0,18\%. Tỷ lệ thấp cho thấy kết quả không đến từ việc thường xuyên chờ trần 48, nhưng first-commit latency cao hơn đáng kể Fixed-N/8 (8 token; 2.687 ms).

\section{Hạn chế và hướng phát triển}
Nguồn thời gian là mô phỏng, không có audio, streaming ASR hoặc lỗi nhận dạng; AL/LAAL token-level không thay thế latency wall-clock. Corpus chỉ có 17 talk TED trong domain kỹ thuật, một provider và một hướng EN--VI; frozen DEV chỉ gồm ba talk nên ước lượng nhạy theo talk. TEST held-out chưa được chạy, do đó không thể claim held-out generalization. Thí nghiệm chỉ dùng một frozen translator, nên kết luận về history chưa chắc chuyển sang model khác. Policy là logistic regression nhẹ; history text chỉ gồm unit vừa commit, không khảo sát history dài hay neural policy lớn. Translator không được fine-tune và prefix translation có thể không phải chế độ mà model được tối ưu riêng.

Hướng tiếp theo hợp lý là (i) đánh giá trên luồng ASR/audio với timestamp quan sát thật; (ii) mở rộng talk, domain và translator nhưng giữ protocol paired; (iii) khảo sát history nhiều unit hoặc policy neural có kiểm soát capacity; và (iv) báo confidence interval theo talk hay bootstrap sau khi có TEST hợp lệ. Các mở rộng này cần preregistration mới, không được thay đổi ngầm Dataset v1 theo kết quả hiện tại.

\section{Kết luận}
Bài báo xây dựng một thí nghiệm có kiểm soát để hỏi liệu committed streaming history cải thiện quyết định \listen/\commit{} trong dịch gia tăng Anh--Việt hay không. P0/P1/P2 là history ablation dùng chung EnViT5 cố định, không phải ba translator. Frozen DEV chọn P1 tại threshold 0,60 với 27,09 BLEU và 61,40 chrF2, nhưng LA-2 đạt chất lượng cao hơn và latency thấp hơn theo cả AL lẫn LAAL. So sánh P0/P1/P2 trên các threshold cho kết quả pha trộn: source history và system-generated target history thay đổi điểm vận hành nhưng chưa chứng minh cải thiện quality--latency nhất quán. Fixed-N/8 vẫn là baseline đơn giản mạnh. Kết luận này chỉ áp dụng cho frozen development evaluation; TEST held-out chưa chạy nên chưa có bằng chứng về khả năng khái quát hóa held-out.

\begin{thebibliography}{00}
\bibitem{ma2019stacl}
M. Ma, L. Huang, H. Xiong, R. Zheng, K. Liu, B. Zheng, C. Zhang, Z. He, H. Liu, X. Li, H. Wu, and H. Wang, ``STACL: Simultaneous Translation with Implicit Anticipation and Controllable Latency using Prefix-to-Prefix Framework,'' in \emph{Proc. 57th Annual Meeting of the ACL}, 2019, pp. 3025--3036, doi: 10.18653/v1/P19-1289.

\bibitem{liu2020partial}
D. Liu, G. Spanakis, and J. Niehues, ``Low-Latency Sequence-to-Sequence Speech Recognition and Translation by Partial Hypothesis Selection,'' in \emph{Proc. Interspeech}, 2020, pp. 3620--3624, doi: 10.21437/Interspeech.2020-2897.

\bibitem{zhang2020adaptive}
R. Zhang, C. Zhang, Z. He, H. Wu, and H. Wang, ``Learning Adaptive Segmentation Policy for Simultaneous Translation,'' in \emph{Proc. EMNLP}, 2020, pp. 2280--2289, doi: 10.18653/v1/2020.emnlp-main.178.

\bibitem{papi2022laal}
S. Papi, M. Gaido, M. Negri, and M. Turchi, ``Over-Generation Cannot Be Rewarded: Length-Adaptive Average Lagging for Simultaneous Speech Translation,'' in \emph{Proc. Third Workshop on Automatic Simultaneous Translation}, 2022, pp. 12--17, doi: 10.18653/v1/2022.autosimtrans-1.2.

\bibitem{post2018sacrebleu}
M. Post, ``A Call for Clarity in Reporting BLEU Scores,'' in \emph{Proc. Third Conference on Machine Translation}, 2018, pp. 186--191, doi: 10.18653/v1/W18-6319.

\bibitem{ngo2022mtet}
C. Ngo, T. H. Trinh, L. Phan, H. Tran, T. Dang, H. Nguyen, M. Nguyen, and M.-T. Luong, ``MTet: Multi-domain Translation for English and Vietnamese,'' arXiv:2210.05610, 2022, doi: 10.48550/arXiv.2210.05610.
\end{thebibliography}

\end{document}
